\reading{CR-7}{Credit Scoring and Retail Credit Risk Management}{DEFER}{2}

\begin{crkeybox}[Learning objectives for this reading]
\begin{enumerate}[label=\textbf{\alph*.},leftmargin=2em]
\item Analyze the credit risks and other risks generated by retail banking.
\item Explain the differences between retail credit risk and corporate credit risk.
\item Discuss the ``dark side'' of retail credit risk and the measures that attempt to address the problem.
\item Define and describe credit risk scoring model types, key variables, and applications.
\item Discuss the key variables in a mortgage credit assessment and describe the use of cutoff scores, default rates, and loss rates in a credit scoring model.
\item Discuss the measurement and monitoring of a scorecard performance including the use of cumulative accuracy profile (CAP) and the accuracy ratio (AR) techniques.
\item Describe the customer relationship cycle and discuss the trade-off between creditworthiness and profitability.
\item Discuss the benefits of risk-based pricing of financial services.
\end{enumerate}
\end{crkeybox}

\begin{crtrapbox}[Single-layer reading]
This is the first reading in the volume that exists in only one of the two source
layers. There is no Class Lecture treatment of CR-7, so the body below is the
Crash layer alone, with gap-fills from the GARP curriculum where the objective
asks for more than the Crash layer carries.
\end{crtrapbox}

% =====================================================================
\lo{CR-7 a}{Analyze the credit risks and other risks generated by retail banking.}

Retail banking, which deals with consumer and small business loans (like
mortgages and credit lines), faces various risks. \term{Credit risk}, the
likelihood of borrower default, is the primary concern. Beyond it:

\begin{enumerate}[label=\alph*)]
\item \term{Operational risks.} Arising from daily business operations.
\item \term{Business risks.} Strategies for new products or market trends, and
changes in loan volumes.
\item \term{Reputation risks.} Damage to the bank's image with customers or
regulators.
\item \term{Interest rate risks.} Fluctuations in interest rates impacting the
bank's profitability.
\item \term{Asset valuation risk.} Difficulty in accurately valuing assets (loans,
bonds) due to market changes.
\end{enumerate}

% =====================================================================
\lo{CR-7 b}{Explain the differences between retail credit risk and corporate credit risk.}

\begin{center}
\footnotesize
\begin{tabular}{L{2.9cm} L{5.6cm} L{6.0cm}}
\toprule
 & \thead{Retail credit risk} & \thead{Corporate credit risk} \\
\midrule
Loans & Small loans, to individuals and retail customers &
  Large loans, to businesses and corporations \\
Portfolio diversification & Very well diversified, due to the large number of
  borrowers &
  Poorly diversified and more concentrated, because there are fewer borrowers \\
Predictability & More predictable, due to diversification and credit scoring &
  Less predictable, because it turns on a specific company's performance \\
Early warning signs & Changes in customer behaviour can signal potential defaults
  (e.g., missed payments) &
  May not show clear signs until it is too late (e.g., internal company issues) \\
Management of risk & Easier to manage, due to diversification and early warning
  indicators &
  More challenging to manage, due to larger loans and fewer early warning
  indicators \\
Impact of defaults & Less severe; a single default has little impact &
  Potentially severe, due to large losses \\
\bottomrule
\end{tabular}
\end{center}
\figcap{Every row runs the same way: the feature that makes retail credit easier
to model is the sheer number of small, independent borrowers, and the feature
that makes corporate credit harder is that a few large exposures carry the
portfolio.}

% =====================================================================
\lo{CR-7 c}{Discuss the ``dark side'' of retail credit risk and the measures that attempt to address the problem.}

\begin{crdefbox}[The dark side]
Unexpected events can cause significant losses in retail portfolios. Causes
include:
\begin{enumerate}[label=\alph*)]
\item lack of historical data on new products;
\item economic downturns causing unexpected defaults;
\item evolving regulations or legal systems impacting defaults; and
\item operational flaws in credit assessment processes.
\end{enumerate}
\end{crdefbox}

\subsection*{Consumer Financial Protection Act (CFPA)}

The underwriting standards a lender is asked to consider are:

\begin{itemize}
\item debt-to-income ratio;
\item upfront fees and points;
\item loan type (balloon payments, interest-only, and so on); and
\item loan term length.
\end{itemize}

\gapbadge

{\footnotesize\color{FRMGrey}The source notes list the four underwriting standards
but do not say what the CFPA is, where it came from, or what it requires. The
following is written from the GARP curriculum.}

\begin{crgapbox}[What the objective asks for: the measures that address the problem]
Since the crisis, various industry reforms and regulations, such as the
\term{Consumer Financial Protection Bureau (CFPB)}, have evolved out of the
\term{Dodd-Frank Act (DFA)} to help deal with the dark side of retail credit risk.

\smallskip
The CFPA requires originators of credit to determine whether the consumer has the
\term{ability to repay} the mortgage. If a mortgage is labelled a \term{qualified
mortgage (QM)}, then a creditor can assume the borrower has met this requirement.
The CFPA also introduced an ``ability to repay'' consideration that asks the
lender to consider the underwriting standards listed above.
\end{crgapbox}

% =====================================================================
\lo{CR-7 d}{Define and describe credit risk scoring model types, key variables, and applications.}

Credit scores use applicant information to predict the chance of loan repayment.
A higher score means lower risk for lenders and better rates for borrowers. There
are three main credit scoring models.

\begin{center}
\footnotesize
\begin{tabular}{L{2.5cm} L{5.2cm} L{6.8cm}}
\toprule
\thead{Model} & \thead{Application and character} & \thead{Key variables} \\
\midrule
\term{Credit bureau scores} & General consumer credit assessment. Utilize FICO
  scores ranging from 300 to 850. Quick, cost-effective, and easy to implement &
  Payment history; credit utilization ratio; length of credit history \\
\term{Pooled model} & Credit assessment for specific industries. Built by
  external parties. More expensive, but can be customized for specific industries &
  Payment history; credit utilization ratio; length of credit history;
  \textbf{industry-specific variables} \\
\term{Custom model} & Credit assessment for unique lending products. Developed
  in-house using the lender's own data, allowing personalized assessment of
  applicants &
  Payment history; credit utilization ratio; length of credit history;
  \textbf{credit mix}; \textbf{credit inquiries} \\
\bottomrule
\end{tabular}
\end{center}
\figcap{All three share the same three base variables. What separates them is
what gets added and who builds it: the pooled model adds industry variables and is
built externally; the custom model adds credit mix and credit inquiries and is
built in-house.}

% =====================================================================
\lo{CR-7 e}{Discuss the key variables in a mortgage credit assessment and describe the use of cutoff scores, default rates, and loss rates in a credit scoring model.}

\subsection*{Key variables in a mortgage credit assessment}

\begin{enumerate}[label=\alph*)]
\item \term{FICO score.} Measures default risk based on credit history.
\item \term{Loan-to-value (LTV) ratio.} Mortgage amount divided by the property's
appraised value.
\item \term{Debt-to-income (DTI) ratio.} Monthly debt payments compared to
monthly gross income.
\item \term{Payment types.} Type of mortgage (e.g., adjustable rate, fixed rate).
\item \term{Documentation types.} Full doc; stated income; no ratio; no income /
no asset.
\end{enumerate}

\subsection*{Cutoff scores, default rates and loss rates}

\begin{crdefbox}[Cutoff score]
Cutoff scores are thresholds set by lenders which dictate whether credit will or
will not be extended, as well as the terms associated with the loans.
\end{crdefbox}

Probability of default and loss given default metrics are critical to assessing
the risk associated with various lender portfolios.

% =====================================================================
\lo{CR-7 f}{Discuss the measurement and monitoring of a scorecard performance including the use of cumulative accuracy profile (CAP) and the accuracy ratio (AR) techniques.}

In assessing the performance of a scorecard, a \term{cumulative accuracy profile
(CAP)} and the \term{accuracy ratio (AR)} are often used.

On the horizontal axis is the population sorted by score, from the highest risk
score to the lowest risk score. On the vertical axis is the percentage of actual
defaults, taken from the bank's records.

\begin{enumerate}[label=\alph*)]
\item The \term{perfect model} line represents an ideal scenario where the model
perfectly predicts defaults. If the model were perfect and the default rate is
5\%, all actual defaults would fall in the first 5\% of the score distribution.
\item The \term{random model} line, the 45-degree diagonal, represents no
predictive ability. It cannot differentiate between good and bad customers and is
equivalent to random guessing.
\item The \term{observed model} is the actual performance of the bank's credit
scoring model, and the bank hopes it sits close to the perfect model line.
\end{enumerate}

\begin{center}
\begin{tikzpicture}
\begin{axis}[
  width=11.6cm, height=6.6cm,
  axis lines=left,
  xlabel={Population distribution based on credit scores},
  ylabel={Percentage of actual defaults},
  xlabel style={font=\sffamily\scriptsize}, ylabel style={font=\sffamily\scriptsize},
  tick label style={font=\sffamily\scriptsize},
  xmin=0, xmax=1, ymin=0, ymax=1.16,
  xtick={0,0.05,0.25,0.5,0.75,1}, xticklabels={0\%,5\%,,50\%,,100\%},
  ytick={0,0.5,1}, yticklabels={0\%,50\%,100\%},
  clip=true,
  legend style={font=\scriptsize\sffamily, draw=FRMRule, fill=white,
    at={(0.5,-0.30)}, anchor=north, legend columns=-1, column sep=8pt}]
% perfect
\addplot[draw=FRMGreen, line width=1.1pt, sharp plot, name path=perf]
  coordinates {(0,0) (0.05,1) (1,1)};
\addlegendentry{perfect model (5\% default rate)}
% observed
\addplot[domain=0.0001:1, samples=180, smooth, draw=FRMNavy, line width=1.2pt,
  name path=obs] {x^0.35};
\addlegendentry{observed model}
% random
\addplot[draw=FRMGrey, line width=1.0pt, dashed, sharp plot, name path=rand]
  coordinates {(0,0) (1,1)};
\addlegendentry{random model}
\addplot[FRMGold!45] fill between[of=obs and rand];
\addplot[FRMRed!22] fill between[of=perf and obs];
\node[font=\sffamily\small\bfseries, text=FRMGold, fill=white, inner sep=1.5pt]
  at (axis cs:0.60,0.62) {$A_R$};
\node[font=\sffamily\scriptsize\bfseries, text=FRMRed, fill=white, inner sep=1.5pt]
  at (axis cs:0.28,0.93) {the rest of $A_P$};
\end{axis}
\end{tikzpicture}
\end{center}
\figcap{The cumulative accuracy profile. $A_R$ is the area between the observed
curve and the random line; $A_P$ is the whole area between the perfect line and
the random line, that is, the gold region plus the pink one. A model with no
predictive power lies \emph{on} the diagonal, so its $A_R$ is zero. Note the
direction of the horizontal axis: the highest-risk scores sit at the left-hand
end, which is why a good model climbs steeply at the start.}

\begin{crfmlbox}[Accuracy ratio]
\[ \mathrm{AR} \;=\; \frac{A_R}{A_P} \]
\vk{$A_R$ = the area captured by the actual rating model, measured against the
random line; $A_P$ = the corresponding area for the perfect model. The closer
this ratio is to 1, the more accurate the model.}
\end{crfmlbox}

\begin{crtrapbox}[Correction: what a random model scores]
The source notes instruct the reader to remember ``Ratio $=1$, model is perfect;
Ratio $=0.5$, model is random.'' The first half is right and the second half is
not, and the notes contradict themselves on the same line: a random model sits on
the 45-degree line, so $A_R = 0$, so $\mathrm{AR} = 0/A_P = \mathbf{0}$.

\smallskip
\textbf{The accuracy ratio runs from 0 (random) to 1 (perfect).} The value 0.5 is
the random benchmark for a \emph{different} statistic, the area under the ROC
curve (AUROC), which runs from 0.5 to 1. The two are related by
$\mathrm{AR} = 2\times\mathrm{AUROC} - 1$.
\end{crtrapbox}

The performance of a scoring model can be monitored, say every quarter, by means
of a CAP curve, and the model replaced when its performance deteriorates.

% =====================================================================
\lo{CR-7 g}{Describe the customer relationship cycle and discuss the trade-off between creditworthiness and profitability.}

Beyond just creditworthiness, lenders also consider potential profitability when
issuing loans. They use various scorecards to assess both credit risk and
potential customer value.

The customer relationship cycle involves:

\begin{enumerate}
\item \term{Marketing} products;
\item \term{Screening} applicants, using scorecards to decide who gets approved
and at what price;
\item \term{Managing accounts}, including pricing and credit line management; and
\item \term{Cross-selling} additional products to existing customers.
\end{enumerate}

% =====================================================================
\lo{CR-7 h}{Discuss the benefits of risk-based pricing of financial services.}

\begin{crdefbox}[Risk-based pricing]
Risk-based pricing (RBP) involves charging different prices for the same product
such that higher (lower) prices can be charged to higher (lower) risk customers.
\end{crdefbox}

Several external and internal factors are used to determine the prices charged,
which are then evaluated in conjunction with various key performance metrics at
each score (risk) band, in order to maximize the trade-off between risk and
profitability.
